\section{领域二：工具设计与 MCP 集成（18\%）}

本领域聚焦于如何设计高质量的工具描述（tool description）、理解 MCP 协议，以及正确配置 \texttt{tool\_choice} 参数。

\subsection{工具描述是工具选择的核心}

\begin{importantbox}{工具描述决定 Claude 的工具选择行为}
Tool description 是 Claude 用于决定调用哪个工具的\textbf{主要依据}。如果描述模糊或多个工具描述重叠，工具选择将变得不可靠。
\end{importantbox}

原文举了一个典型案例：\texttt{get\_customer} 和 \texttt{lookup\_order} 两个工具的描述过于相似，导致 Claude 经常选错工具。面对这种情况，正确的修复方案不是：
\begin{itemize}[leftmargin=1.5em]
    \item 添加 few-shot 示例（治标不治本）
    \item 构建路由分类器（过度工程）
    \item 合并工具（丧失功能边界）
\end{itemize}
\textbf{正确做法}：重写工具描述，使每个工具的职责边界清晰、不重叠。

\subsection{tool\_choice 参数详解}

\begin{table}[H]
\centering
\begin{tabular}{lp{9cm}}
\toprule
\textbf{选项} & \textbf{行为} \\
\midrule
\texttt{"auto"} & 模型自行决定是调用工具还是直接回复文本 \\
\texttt{"any"} & 模型\textbf{必须}调用某个工具，但由模型选择具体调用哪一个 \\
\texttt{forced selection} & 模型\textbf{必须}调用指定的特定工具 \\
\bottomrule
\end{tabular}
\caption{\texttt{tool\_choice} 三种模式对比}
\end{table}

\subsection{工具数量的上限原则}

\begin{warningbox}{单个 Agent 的工具不应超过 4--5 个}
给一个 Agent 提供 18 个工具会严重降低工具选择的准确率。正确做法是将工具按角色分配给不同的子 Agent，每个子 Agent 只携带 4--5 个与其职责相关的工具。
\end{warningbox}

\subsection{推荐实践项目}

\begin{knowledgebox}{动手练习：工具描述对比实验}
创建两个功能相似的 MCP 工具，故意写模糊的描述使其产生路由错误。然后修复描述，亲身体验描述质量对工具选择的影响。
\end{knowledgebox}

\subsection{本章小结}

工具设计的关键在于：写出清晰、无歧义的工具描述；熟练掌握 \texttt{tool\_choice} 的三种模式及适用场景；控制单个 Agent 的工具数量在 4--5 个以内，通过子 Agent 分工来扩展能力。

